% SW403 Course Project -- Phase 1: Proposal & Background Study
% Compile on Overleaf (pdfLaTeX). Target length: 2-3 pages.
\documentclass[conference]{IEEEtran}
\usepackage{cite}
\usepackage{amsmath}
\usepackage{booktabs}
\usepackage{array}
\usepackage{url}

\begin{document}

\title{Does Multi-Agent Collaboration Really Help LLM Bug Fixing?\\
A Budget-Controlled, Contamination-Aware Study}

% TODO: replace with real names, IDs and emails (2-3 members).
\author{
\IEEEauthorblockN{Member 1 (ID)}
\IEEEauthorblockA{Zewail City of Science and Technology\\ email}
\and
\IEEEauthorblockN{Member 2 (ID)}
\IEEEauthorblockA{Zewail City of Science and Technology\\ email}
\and
\IEEEauthorblockN{Member 3 (ID)}
\IEEEauthorblockA{Zewail City of Science and Technology\\ email}
}

\maketitle

\begin{abstract}
LLM-based automated program repair (APR) increasingly relies on multi-agent
pipelines in which specialised agents localise, patch and review bugs. Reported
gains are hard to interpret for three reasons: multi-agent systems usually spend
more LLM calls than the baselines they are compared with; patches that pass the
available tests (plausible) are often still incorrect; and popular benchmarks
such as QuixBugs are public, likely part of LLM training data, and already
saturated. We propose a controlled study and prototype that compares a
single-shot LLM, a single-agent iterative repair loop and a LangGraph
multi-agent pipeline (localiser, fixer, test-runner tool, reviewer) under an
identical per-bug LLM-call budget. We evaluate on QuixBugs-Python and a sample
of the leakage-oriented ConDefects-Python dataset, measuring plausible and
correct fix rates on held-out tests, overfitting rate, a hallucination taxonomy,
cost, and a memorisation probe based on identifier-renamed programs. The
prototype is exposed as a FastAPI service with a Streamlit dashboard.
\end{abstract}

\begin{IEEEkeywords}
automated program repair, large language models, multi-agent systems,
hallucination, data contamination
\end{IEEEkeywords}

\section{Introduction and Problem Definition}
Automated program repair has been studied for over a decade, from search-based
tools such as GenProg~\cite{genprog} to today's LLM-based systems. LLMs changed
the field quickly: conversational repair with test feedback fixed 162 of 337
Defects4J bugs at about \$0.42 per bug~\cite{chatrepair}; autonomous agents that
call tools repaired 164 Defects4J bugs~\cite{repairagent}; and a hierarchical
multi-agent framework reports fixing all QuixBugs programs~\cite{unidebugger}.

These results are difficult to compare, and three confounds stand out.
\textbf{(1) Budget.} A multi-agent pipeline usually makes more LLM calls than
its baseline, so part of its advantage may simply be extra compute. Agentless, a
fixed localise-then-repair pipeline without an agent loop, was competitive with
agent-based systems on SWE-bench Lite~\cite{agentless}.
\textbf{(2) Plausible is not correct.} Qi et al.\ showed that the overwhelming
majority of patches reported by earlier generate-and-validate tools were not
correct~\cite{qi2015}; a passing test suite is weak evidence of a fix.
\textbf{(3) Contamination and saturation.} QuixBugs~\cite{quixbugs} dates from
2017 and its fixes are public. ChatGPT fixed 31/40 of its Python bugs with
dialogue~\cite{sobania}, GPT-4o 38/40 and o1 40/40~\cite{gpto1}. A benchmark at
its ceiling cannot show improvement, and memorisation cannot be separated from
repair ability. A related concern is that LLMs struggle to self-correct
without external feedback~\cite{selfcorrect}, so a ``reviewer'' agent may
approve wrong patches rather than catch them.

\textbf{Problem statement.} Does multi-agent collaboration improve LLM-based bug
fixing once budget, patch correctness and contamination are controlled?

\textbf{Planned contributions.} (a) A budget-matched comparison of single-shot,
single-agent and multi-agent repair with the same model; (b) a held-out-test
correctness check and a hallucination taxonomy for repair agents; (c) a
memorisation probe using semantics-preserving identifier renaming; (d) an open
prototype (FastAPI, LangGraph, Streamlit).

\section{Related Work}
\textbf{Patch correctness.} Qi et al.~\cite{qi2015} distinguish plausible
patches (pass the validation tests) from correct ones and find that most
reported patches of GenProg, RSRepair and AE were not correct, many equivalent
to deleting functionality. We adopt this distinction as a core metric.

\textbf{LLM-based repair.} Sobania et al.~\cite{sobania} evaluated ChatGPT on
QuixBugs-Python and found that follow-up hints raised its fixes to 31/40.
ChatRepair~\cite{chatrepair} interleaves patch generation with test-failure
feedback in a conversation. Self-Debugging~\cite{selfdebug} shows that execution
feedback helps LLMs correct their own code. These are single-agent loops; our
single-agent baseline follows this design.

\textbf{Agents and multi-agent systems.} SWE-agent~\cite{sweagent} and
RepairAgent~\cite{repairagent} let one LLM plan and call tools autonomously.
UniDebugger (earlier FixAgent)~\cite{unidebugger} organises several agents
hierarchically for localisation and repair. Agentless~\cite{agentless} is a
counterpoint: a fixed pipeline without an agent loop was competitive on
SWE-bench Lite at low cost.

\textbf{Benchmarks.} Defects4J~\cite{defects4j} (Java) and SWE-bench~\cite{swebench}
(repository-level Python) are standard but expensive to run in a course
setting. HumanEvalFix~\cite{octopack} provides 164 human-written buggy
functions with tests. ConDefects~\cite{condefects} collects 1,625 Python and
1,254 Java faulty AtCoder programs written between October 2021 and September
2023, specifically to reduce training-data leakage.

\begin{table}[t]
\caption{Representative prior work}
\label{tab:rw}
\centering
\footnotesize
\begin{tabular}{@{}>{\raggedright}p{1.35cm}>{\raggedright}p{1.6cm}>{\raggedright}p{1.35cm}>{\raggedright\arraybackslash}p{2.9cm}@{}}
\toprule
Work & Approach & Benchmark & Key result / limitation for us \\
\midrule
Sobania~\cite{sobania} & ChatGPT dialogue & QuixBugs-Py & 31/40 with hints; benchmark likely in training data \\
ChatRepair~\cite{chatrepair} & Single agent, test feedback & Defects4J & 162/337 at \$0.42/bug; single agent only \\
RepairAgent~\cite{repairagent} & Autonomous tool-using agent & Defects4J (835) & 164 fixed; one agent, not multi-agent \\
UniDebugger~\cite{unidebugger} & Hierarchical multi-agent & QuixBugs + others & Reports all QuixBugs fixed; benchmark at ceiling \\
Agentless~\cite{agentless} & Fixed pipeline, no agent & SWE-bench Lite & Competitive with agents; questions need for agents \\
Hu et al.~\cite{gpto1} & GPT family comparison & QuixBugs-Py & o1 40/40, GPT-4o 38/40; saturation \\
\bottomrule
\end{tabular}
\end{table}

\section{Research Gap and Questions}
Table~\ref{tab:rw} shows that multi-agent repair is mostly reported as fix
counts on benchmarks that are saturated and likely contaminated, and we are not
aware of a study that jointly controls LLM-call budget, patch overfitting and
contamination at a scale a small team can replicate. We ask:

\textbf{RQ1 (Effectiveness under equal budget).} Does a multi-agent pipeline fix
more bugs than a single-agent iterative loop when both receive the same per-bug
LLM-call budget? \emph{Metrics:} plausible and correct fix rate, paired per bug.

\textbf{RQ2 (Validation and hallucination).} Does a reviewer agent reduce
plausible-but-incorrect patches and hallucinations, or does it approve them?
\emph{Metrics:} overfitting rate, reviewer precision/recall against held-out
test outcomes, hallucination counts per category.

\textbf{RQ3 (Contamination).} How much does repair success drop on
semantics-preserving, identifier-renamed versions of the same programs, and
does the drop differ between QuixBugs and ConDefects?
\emph{Metric:} change in correct fix rate.

\section{Proposed Approach and Evaluation Plan}
\textbf{Datasets.} (i) QuixBugs-Python: 40 programs, each with a one-line bug,
test cases and a reference fix~\cite{quixbugs}. (ii) A random sample of about
100 ConDefects-Python bugs stratified by difficulty, using the dataset's
test cases~\cite{condefects}. (iii) Renamed variants of (i) and (ii), produced
with Python's \texttt{ast} module and checked by running the reference fix.
For each bug, tests are split into \emph{visible} tests (given to the
systems) and \emph{held-out} tests (used only for evaluation).

\textbf{Systems.} All systems use the same LLM, temperature and a budget of
$K$ LLM calls per bug (e.g., $K=8$); token usage is logged because calls differ
in length.
\begin{itemize}
\item \textbf{S0, single-shot:} one prompt, one patch.
\item \textbf{S1, single-agent loop:} patch, run visible tests, return failures
to the model, repeat until pass or budget exhausted~\cite{selfdebug,chatrepair}.
\item \textbf{S2, multi-agent (LangGraph):} Localiser (LLM) $\rightarrow$ Fixer
(LLM) $\rightarrow$ Test-runner (deterministic sandboxed tool, not an LLM)
$\rightarrow$ Reviewer (LLM, approve or reject with reason), looping back to the
Fixer on failure or rejection.
\item \textbf{S2$-$R:} S2 without the Reviewer (ablation for RQ2).
\end{itemize}
We plan one commercial API model and one open-weights code model (e.g.,
Qwen2.5-Coder) to test whether findings generalise, subject to cost.

\textbf{Evaluation module (automatic checker).} Applies each patch in a
sandboxed subprocess with a timeout, runs visible and held-out tests and labels
the outcome \emph{correct}, \emph{plausible-but-incorrect} or \emph{failing}.
It also flags hallucinations: (H1) non-runnable code; (H2) non-existent API or
import; (H3) edits outside the buggy function or to the tests; (H4) unchanged
code reported as fixed; (H5) reviewer approval of a patch that fails held-out
tests. Two members manually compare a random 20\% of plausible patches with the
reference fix and report Cohen's $\kappa$.

\textbf{Analysis.} Three runs per bug and condition; we report means with 95\%
confidence intervals and use McNemar's test on paired per-bug correctness
(S1 vs.\ S2, S2 vs.\ S2$-$R).

\textbf{Prototype.} FastAPI endpoint \texttt{POST /repair} accepting
\{code, tests, mode\} and returning \{patch, verdict, agent trace, tokens\};
LangGraph orchestration; a Streamlit dashboard for fix rates, hallucination
breakdown, cost and per-bug traces.

\textbf{Threats and expected SOTA answer.} Contamination cannot be removed
entirely: current models may have been trained after ConDefects' time window,
which RQ3 partly measures. QuixBugs has only 40 bugs, so intervals will be
wide. Since o1 and UniDebugger already report 40/40 on QuixBugs, we do not
expect to beat the state of the art in raw fix count there; our contribution is
the controlled comparison and the correctness and contamination analysis.

\section{Timeline and Work Split}
\textbf{Wk 7:} proposal. \textbf{Wk 8:} dataset loaders, sandbox, S0/S1.
\textbf{Wk 9:} S2 in LangGraph, evaluation module, FastAPI.
\textbf{Wk 10:} early results on QuixBugs (Phase 2).
\textbf{Wk 11--12:} ConDefects sample, renaming probe, manual labelling.
\textbf{Wk 13:} dashboard and statistics. \textbf{Wk 14:} paper and demo.
% TODO: assign names.
\emph{Member 1:} datasets and evaluation module. \emph{Member 2:} agents and
FastAPI. \emph{Member 3:} dashboard, manual labelling and statistics. All members
write the report and commit their own work.

\section*{AI Tools Declaration}
% TODO: update to reflect what the team actually used.
Claude Code (Anthropic) was used to search for and verify literature and to
draft this proposal. All content was reviewed and edited by the authors.

\begin{thebibliography}{99}
\bibitem{genprog}
C.~Le~Goues, T.~Nguyen, S.~Forrest, and W.~Weimer, ``GenProg: A generic method
for automatic software repair,'' \emph{IEEE Trans. Softw. Eng.}, vol.~38,
no.~1, pp.~54--72, 2012.

\bibitem{qi2015}
Z.~Qi, F.~Long, S.~Achour, and M.~Rinard, ``An analysis of patch plausibility
and correctness for generate-and-validate patch generation systems,'' in
\emph{Proc. ISSTA}, 2015.

\bibitem{quixbugs}
D.~Lin, J.~Koppel, A.~Chen, and A.~Solar-Lezama, ``QuixBugs: A multi-lingual
program repair benchmark set based on the Quixey challenge,'' in \emph{Proc.
SPLASH Companion}, 2017.

\bibitem{defects4j}
R.~Just, D.~Jalali, and M.~D.~Ernst, ``Defects4J: A database of existing faults
to enable controlled testing studies for Java programs,'' in \emph{Proc.
ISSTA}, 2014.

\bibitem{sobania}
D.~Sobania, M.~Briesch, C.~Hanna, and J.~Petke, ``An analysis of the automatic
bug fixing performance of ChatGPT,'' in \emph{Proc. Int. Workshop on Automated
Program Repair (APR)}, 2023, pp.~23--30.

\bibitem{chatrepair}
C.~S.~Xia and L.~Zhang, ``Automated program repair via conversation: Fixing 162
out of 337 bugs for \$0.42 each using ChatGPT,'' in \emph{Proc. ISSTA}, 2024.

\bibitem{selfdebug}
X.~Chen, M.~Lin, N.~Sch\"arli, and D.~Zhou, ``Teaching large language models to
self-debug,'' in \emph{Proc. ICLR}, 2024.

\bibitem{repairagent}
I.~Bouzenia, P.~Devanbu, and M.~Pradel, ``RepairAgent: An autonomous, LLM-based
agent for program repair,'' in \emph{Proc. ICSE}, 2025.

\bibitem{sweagent}
J.~Yang, C.~E.~Jimenez, A.~Wettig, K.~Lieret, S.~Yao, K.~Narasimhan, and
O.~Press, ``SWE-agent: Agent-computer interfaces enable automated software
engineering,'' in \emph{Proc. NeurIPS}, 2024.

\bibitem{agentless}
C.~S.~Xia, Y.~Deng, S.~Dunn, and L.~Zhang, ``Agentless: Demystifying LLM-based
software engineering agents,'' in \emph{Proc. FSE}, 2025, arXiv:2407.01489.

\bibitem{unidebugger}
C.~Lee \emph{et al.}, ``UniDebugger: Hierarchical multi-agent framework for
unified software debugging,'' in \emph{Proc. EMNLP}, 2025, arXiv:2404.17153.

\bibitem{gpto1}
H.~Hu, Y.~Shang, G.~Xu, C.~He, and Q.~Zhang, ``Can GPT-O1 kill all bugs? An
evaluation of GPT-family LLMs on QuixBugs,'' in \emph{Proc. Int. Workshop on
Automated Program Repair (APR)}, 2025, arXiv:2409.10033.

\bibitem{selfcorrect}
J.~Huang \emph{et al.}, ``Large language models cannot self-correct reasoning
yet,'' in \emph{Proc. ICLR}, 2024.

\bibitem{swebench}
C.~E.~Jimenez \emph{et al.}, ``SWE-bench: Can language models resolve
real-world GitHub issues?'' in \emph{Proc. ICLR}, 2024.

\bibitem{octopack}
N.~Muennighoff \emph{et al.}, ``OctoPack: Instruction tuning code large
language models,'' in \emph{Proc. ICLR}, 2024.

\bibitem{condefects}
Y.~Wu \emph{et al.}, ``ConDefects: A complementary dataset to address the data
leakage concern for LLM-based fault localization and program repair,'' in
\emph{Proc. FSE Companion (Demo)}, 2024, arXiv:2310.16253.
\end{thebibliography}

\end{document}
